\documentclass[11pt,a4paper]{article}
\usepackage[margin=0.85in]{geometry}
\usepackage{booktabs}
\usepackage{array}
\usepackage{multirow}
\usepackage[table]{xcolor}
\definecolor{mismatch}{RGB}{250,225,225}
\definecolor{theirs}{RGB}{243,243,243}
\definecolor{ours}{RGB}{225,240,250}
\definecolor{pend}{RGB}{255,247,224}
\usepackage{amssymb,amsmath}
\newcolumntype{L}[1]{>{\raggedright\arraybackslash}p{#1}}
\newcolumntype{C}[1]{>{\centering\arraybackslash}p{#1}}
\setlength{\emergencystretch}{3em}
\usepackage[htt]{hyphenat}
\title{HARP: released code vs.\ published paper\\[2pt]
\large The setup mismatch, and what it is worth in AUROC}
\author{}
\date{\today}
\begin{document}
\maketitle

\noindent\textbf{Purpose.} We reran HARP's released code on our own generated data so that our
comparison against their published numbers is like-for-like. While reading their code to build the
data adapter we found one place where the released implementation does not match the setup the paper
describes. This document reports the audit, the reproduction, and a controlled experiment that
measures what the mismatch is worth.

\medskip
\noindent\textbf{Summary.} The \emph{method} is implemented exactly as the paper describes --
detector architecture, token aggregation, optimiser, and correctness thresholds all match. The one
substantive difference is \textbf{how the data is split into training and test sets}, and on our data
it is worth \textbf{3.8 to 13.3 AUROC points}.

\section*{1. The audit}

\renewcommand{\arraystretch}{1.3}
\begin{table}[h]
\centering\footnotesize
\begin{tabular}{L{2.2cm} L{4.7cm} L{5.4cm} L{1.6cm}}
\toprule
\textbf{Item} & \textbf{What the paper says} & \textbf{What the code does} & \textbf{Match?} \\
\midrule

\textbf{Known / unknown grouping}
& \S3.1: a question is ``known'' if any of its generated answers matches the reference. Appendix A:
the unknown questions all go to the test set.
& Correct. \texttt{get\_known.py:115} marks a question known if any of its 10 answers is judged
correct; \texttt{main.py:264} puts \emph{every} unknown question into the test set and none into
training; the 75/25 ratio is applied only to the known group.
& Yes \\

\rowcolor{mismatch}
\textbf{How the known group is split}
& \S3.1: ``Let $\mathcal{X}_{known} = \{x \mid known(x)=1\}$ denote \textbf{the set of all inputs}
whose knowledge is known to the LLMs.'' \newline
Appendix A: ``75\% of $\mathcal{X}_{known}$ is used for training, while the remaining 25\%,
together with $\mathcal{X}_{unknown}$, is used to test the hallucination detector on
\textbf{unseen data}.''
& Within the known group, splits \textbf{individual answers}, not questions.
\texttt{main.py:184--201} builds one entry per answer;
\texttt{get\_hidden\_state.py:72--79} stores one vector per answer;
\texttt{main.py:237--246} zips these into one list;
\texttt{main.py:259} splits that list 75/25 with a \texttt{torch.randperm} shuffle
(\texttt{utils.py:49--58}).
& \textbf{No} \\

\textbf{Answers per question}
& Appendix A: ``beam search is used to generate 10 answer paths per question.''
& 10 requested (\texttt{main.py:157--158}), but answers decoding to empty text are dropped
(\texttt{get\_known.py:93--94}), so a question can end up with fewer.
& Mostly \\

\textbf{Detector architecture}
& Appendix A: ``a two-layer MLP with hidden dimension 1024 and ReLU activation.''
& \texttt{mlp\_trainer.py:27--31}: \texttt{Linear(input,1024)} $\to$ \texttt{ReLU} $\to$
\texttt{Linear(1024,1)} $\to$ \texttt{Sigmoid}.
& Yes \\

\textbf{Score for a QA pair}
& Eq.~16: ``The maximum score among all tokens is taken as the hallucination score of the QA pair.''
& \texttt{mlp\_trainer.py:56, 192, 213}: \texttt{masked\_scores.max(dim=-1)}, padding masked out.
& Yes \\

\textbf{Training setup}
& Appendix A: 50 epochs, Adam, lr 1e-4, cosine decay, batch size 128, weight decay 3e-4.
& \texttt{mlp\_trainer.py:282} \texttt{Adam(lr, weight\_decay)}; \texttt{mlp\_trainer.py:12}
\texttt{CosineAnnealingLR}; values from \texttt{main.py}; batch size from the README command.
& Yes \\

\textbf{Correctness label}
& \S5.1: correct if ROUGE-L $>0.7$ or semantic similarity $>0.5$.
& \texttt{DatasetJudge.py:106--112}: same thresholds, plus an unmentioned margin condition whose
default of $-1$ (\texttt{DatasetJudge.py:26}) makes it almost always true.
& Yes (see note) \\

\textbf{Subspace dimension}
& \S5.3: ``a dimension of 256 yields the best performance.'' Fixed globally; every Table~1 value is
reported at it.
& \texttt{main.py:59} sweeps $\{32,\dots,1024\}$ and prints all; \texttt{Reasoning\_Patch.ipynb:263}
hardcodes 256.
& Yes \\

\bottomrule
\end{tabular}
\end{table}

\section*{2. The mismatch in plain terms}

\noindent\textbf{What is \emph{not} wrong.} Questions are correctly labelled known or unknown per the
paper's definition; all unknown questions are correctly placed in the test set and never used for
training; the 75/25 ratio is correctly applied only to the known group. There is a genuine
known/unknown separation.

\medskip
\noindent\textbf{What is wrong} is one level further down. Each question is answered ten times. The
paper splits the known group \emph{over questions}; the code pools that group's answers into one list
and splits \emph{that} at random. A known question therefore keeps all ten answers on the training
side only about 5.6\% of the time ($0.75^{10}$) and almost never lands entirely in the test set, so
roughly \textbf{94\% of known questions contribute answers to both sides}. For those, the question is
not unseen at test time.

\medskip
\noindent This is directly measurable, and the run below confirms it: the number of known questions
appearing on both sides of the split was 626 of 666 (TruthfulQA), 5{,}956 of 6{,}316 (TriviaQA), 334
of 360 (NQ-Open) and 278 of 302 (TyDiQA-GP) -- \textbf{92--94\%}, against the 94.4\% predicted.

\section*{3. The reproduction}

We fed our own pinned generations into their generation cache (\texttt{main.py:120--123}) so their
beam search and BLEURT judge never run, and executed their \texttt{main.py} otherwise untouched --
the only edit being two device strings (\texttt{main.py:77--78}) so it runs on one GPU rather than
three. All values at their fixed $\text{proj\_dim}=256$.

\begin{table}[h]
\centering\footnotesize
\begin{tabular}{lccccc}
\toprule
\textbf{Qwen2.5-7B} & \textbf{TruthfulQA} & \textbf{TriviaQA} & \textbf{NQ-Open} & \textbf{TyDiQA-GP}
& \textbf{mean $|\Delta|$} \\
\midrule
HARP published (their data) & 88.1 & 92.8 & 84.0 & 88.4 & -- \\
Their code, our data        & 86.27 & 92.89 & 83.12 & 91.87 & -- \\
\midrule
$\Delta$ & $-1.83$ & $+0.09$ & $-0.88$ & $+3.47$ & 1.57 \\
\bottomrule
\end{tabular}
\caption{Mean signed deviation $+0.21$. Their method reproduces on our data with no systematic bias
in either direction, which is what licenses the controlled experiment below.}
\end{table}

\section*{4. The controlled experiment}

Everything is held fixed -- their hidden states, their lm\_head SVD, their projection, their MLP,
their hyperparameters, the same five seeds -- and \textbf{only the unit of the 75/25 split is
varied}. Their own \texttt{utils.split\_data} is called for the answer-level arm, so that arm is
their code path exactly; it doubles as a control and reproduces what their \texttt{main.py} printed.

\begin{table}[h]
\centering\footnotesize
\renewcommand{\arraystretch}{1.35}
\begin{tabular}{L{2.1cm} L{2.9cm} cccc}
\toprule
\textbf{Method} & \textbf{Split protocol} & \textbf{TruthfulQA} & \textbf{TriviaQA} &
\textbf{NQ-Open} & \textbf{TyDiQA-GP} \\
\midrule
\rowcolor{theirs}
\textbf{HARP} & answer-level \newline \emph{(their code)}
& 85.16 {\scriptsize$\pm$1.23} & 92.50 {\scriptsize$\pm$0.24}
& 81.03 {\scriptsize$\pm$2.41} & 91.75 {\scriptsize$\pm$0.95} \\
\rowcolor{theirs}
\textbf{HARP} & question-level \newline \emph{(their paper)}
& 77.54 {\scriptsize$\pm$1.51} & 88.75 {\scriptsize$\pm$0.33}
& 67.73 {\scriptsize$\pm$2.24} & 79.90 {\scriptsize$\pm$2.87} \\
\midrule
\rowcolor{pend}
\textbf{Ours} & answer-level \newline \emph{(their code)}
& \textit{pending} & \textit{pending} & \textit{pending} & \textit{pending} \\
\rowcolor{ours}
\textbf{Ours} & question-level \newline \emph{(their paper)}
& \textbf{87.8} & \textbf{92.7} & \textbf{79.4} & \textbf{88.3} \\
\midrule
\multicolumn{2}{l}{\emph{cost of the mismatch} (HARP)}
& $-7.62$ & $-3.75$ & $-13.30$ & $-11.85$ \\
\multicolumn{2}{l}{\emph{known questions}}
& 666 & 6{,}316 & 360 & 302 \\
\bottomrule
\end{tabular}
\caption{Pooled AUROC (\%) on Qwen2.5-7B, $\text{proj\_dim}=256$, 5 seeds. The two HARP rows differ
in \emph{nothing} but the split unit. Our rows use our unsupervised feature conditions with an
RF/LR readout; the best condition per dataset is shown.}
\end{table}

\noindent\textbf{Two things follow.}

\medskip
\noindent\textbf{(i) The mismatch is worth 3.8 to 13.3 points.} Every dataset drops substantially
when the split is made over questions instead of answers, mean 9.1 points. Their README states that
running the given command obtains the paper's experimental results, so the published numbers were
produced under the answer-level split, not the protocol \S3.1 describes.

\medskip
\noindent\textbf{(ii) Under the paper's own protocol, our unsupervised conditions beat their
supervised detector on all four datasets} -- by 10.3, 4.0, 11.7 and 8.4 points respectively.
Comparing our question-level numbers against their answer-level numbers, as the literature table
currently does, we lose three of four. The protocol was the entire gap.

\section*{5. A prediction we got wrong, and the corrected mechanism}

An earlier draft of this document predicted the effect would scale with the share of test
\emph{answers} drawn from questions seen in training -- large on TruthfulQA (52\%), negligible on
NQ-Open (3\%). \textbf{The data refutes this.} NQ-Open has the smallest such share and the largest
effect ($-13.30$).

The error was in the quantity. AUROC is built from (hallucinated, correct) pairs, and \emph{every
such pair requires a correct answer}. By construction a correct answer can only occur inside a known
question -- an unknown question is one where all ten answers were judged wrong. Unknown questions
therefore contribute only one class and form no discriminative pairs at all. On NQ-Open the 32{,}500
unknown answers are 97\% of the test set and contribute nothing to the metric; all $\approx$280
correct answers sit in the known group, where 93\% of questions straddle the split. The leakage
touches essentially \emph{all} of the usable signal, on every dataset. That is why the effect is
large everywhere.

What the effect size does track, cleanly and inversely, is the \textbf{number of distinct known
questions}: TyDiQA-GP 302 $\to$ 11.85, NQ-Open 360 $\to$ 13.30, TruthfulQA 666 $\to$ 7.62, TriviaQA
6{,}316 $\to$ 3.76. With few questions to memorise, memorisation is cheap and there is little
diversity to generalise from; with six thousand, the detector is pushed toward a transferable rule.

\section*{6. Caveats}

\begin{itemize}
\item \textbf{The missing cell.} Until our method is also run under the answer-level split, the
comparison in Table~3 varies method and protocol together in the top-vs-bottom direction. The
HARP-vs-HARP contrast is fully controlled; the HARP-vs-ours contrast is not yet.
\item \textbf{Readouts differ.} Our conditions use a random forest or logistic regression on pooled
multi-layer features; theirs is an MLP over per-token projections with a max aggregation. That is
the method difference under test, but it is not an ablation of the split alone across both systems.
\item \textbf{NQ-Open is noisy.} Its question-level arm carries $\pm2.24$ across seeds, because only
$\approx$280 of 33{,}400 validation answers are correct. Treat its individual value loosely; the
direction is unambiguous.
\item \textbf{TriviaQA labels differ slightly.} Their \texttt{DatasetInit.py:103} folds
\texttt{normalized\_aliases} into the correct-answer list and ours does not, so their published 92.8
was computed over a marginally different known/unknown partition. The near-exact reproduction on that
dataset ($+0.09$) suggests the difference is cosmetic, but it is unmeasured.
\item \textbf{Their LLaMA is the base model} (\texttt{main.py:20}, \texttt{meta-llama/Llama-3.1-8B}),
ours is Instruct. Paper and code agree with each other, so this is our deviation, not theirs -- which
is why every number here is Qwen, where the models match exactly.
\end{itemize}

\medskip
\noindent This reads as a mismatch between the write-up and the released implementation rather than
anything deliberate: the paper's formalism is precise, and the method itself is implemented
faithfully. The consequence for the field is the same either way.

\end{document}
